% !PRINTED 590
\weberpage{590}
%<<CONT
which are relatively prime to one another, and one can then
determine two odd integers $\alpha_1, \alpha_2$ in $\Omega$ that are relatively
prime to one another, of which the one is divisible by $\mathfrak{a}_1$, the
other by $\mathfrak{a}_2$. Then $\eta\alpha_1\alpha_2$ is an integer in $\Omega$,
and by Bd.\ II, \S\,166, (4) one can determine two integers $\varrho_1, \varrho_2$
in such a way that
\[
\frac{\nu\,\alpha_1\alpha_2}{\mu} = \varrho_1\alpha_2 + \varrho_2\alpha_1
\]
and therefore
\[
\boldsymbol{s}\!\left(\frac{2\nu}{\mu}\right) = \boldsymbol{s}\!\left(\frac{2\varrho_1}{\alpha_1} + \frac{2\varrho_2}{\alpha_2}\right)
\]
holds. But this is of the form (25). From it we obtain
the theorem:

4. \gesperrt{If a fractional number $\eta$ in $\Omega$ with odd
denominator contains two or more distinct prime factors in the
denominator, then $\boldsymbol{s}(\eta)$ is a unit.}

\begin{center}
\textbf{\S\ 158. \gesperrt{Decomposition of the Function $\boldsymbol{A}(x)$.}}
\end{center}

Now let
\weq{1}{\begin{aligned} \mu &= y + x\sqrt{\Delta}, \\ \mu_1 &= y_1 + x_1\sqrt{\Delta} \end{aligned}}
be two odd integers in $\Omega$ of the form \S\,157, (6). Let
$\mathfrak{m}$ be the greatest common ideal divisor of $\mu$ and $\mu_1$.
If one therefore sets
\weq{2}{\begin{aligned} \mu &= \mathfrak{m}\mathfrak{a}, \\ \mu_1 &= \mathfrak{m}\mathfrak{a}_1 \end{aligned}}
then $\mathfrak{a}$ and $\mathfrak{a}_1$ are two \gesperrt{equivalent ideals} without a common
divisor.

The functions $\boldsymbol{A}(x)$ that correspond to the two numbers $\mu, \mu_1$
we denote by
\[
\boldsymbol{A}_\mu(x), \ \boldsymbol{A}_{\mu_1}(x),
\]
and we ask what common roots these functions
have, that is, when the equation:
\weq{3}{\boldsymbol{s}\!\left(\frac{2\varrho}{\mu}\right) = \boldsymbol{s}\!\left(\frac{2\varrho_1}{\mu_1}\right)}
can be satisfied, when $\varrho$ is taken with respect to the modulus $\mu$, and $\varrho_1$
with respect to the modulus $\mu_1$.
